\documentclass[11pt,a4paper]{article}
\usepackage{amsmath,amsfonts,amssymb}
\usepackage{graphicx}
\usepackage{booktabs}
\usepackage{hyperref}
\usepackage[authoryear]{natbib}
\usepackage{geometry}
\geometry{margin=1in}

\title{Optimización Estratégica en Telecomunicaciones:\\ Un Enfoque Guiado por Rentabilidad empleando Inteligencia Artificial Explicable (XAI) e Inteligencia de Negocios (BI)}
\author{A. Autor \\ Instituto de Investigación \\ \texttt{autor@instituto.edu}}
\date{\today}

\begin{document}

\maketitle

\begin{abstract}
Este documento investiga empíricamente cómo la integración de modelos predictivos de Machine Learning, enriquecidos con Inteligencia Artificial Explicable (XAI) y una matriz de evaluación coste-beneficio, optimiza las decisiones de retención de clientes (Churn) dentro de sistemas operativos de Inteligencia de Negocios (BI). Utilizando datos históricos públicos de clientes en el sector de las telecomunicaciones, demostramos cómo traducir probabilidades puramente predictivas en rentabilidad cuantificable (ROI) mejora las estrategias comerciales focalizadas.
\end{abstract}

\section{Introducción}
La industria de las telecomunicaciones enfrenta uno de los entornos más competitivos y saturados del mercado global. En este contexto, la retención de clientes (generalmente referida como mitigación de \textit{Churn}) ha superado la adquisición de clientes como el principal impulsor de la rentabilidad a largo plazo. 

Las infraestructuras de Machine Learning (ML) se han posicionado fuertemente como herramientas predictivas para anticipar qué usuarios tienen la mayor propensión a cancelar sus servicios. Sin embargo, la adopción operativa de estos algoritmos a menudo fracasa en la última milla por dos razones críticas: 1) la métrica de éxito de ML (usualmente Precisión o AUC) no siempre se alinea con la rentabilidad comercial real (ROI), y 2) los resultados se entregan como cajas negras probabilísticas que no brindan explicaciones accionables a los agentes humanos.

Este documento propone arquitecturar y evaluar de forma empírica una solución integral de toma de decisiones dirigida al valor de negocio (\textit{Profit-driven}). Mediante el cruce de modelos ensamblados (XGBoost), Inteligencia Artificial Explicable (XAI) basado en valores SHAP, y tableros interactivos de Inteligencia de Negocios (BI), buscamos responder a la siguiente pregunta: ¿Cómo la integración visual de explicabilidad algorítmica y matrices de coste-beneficio empodera y rentabiliza las campañas de retención en telecomunicaciones?


\section{Revisión de Literatura}
La literatura previa en análisis predictivo de Churn se ha enfocado primariamente en la maximización del rendimiento algorítmico, favoreciendo ensambles de árboles (por ejemplo, Random Forest o Gradient Boosting) sobre modelos lineales tradicionales dada su destreza en esquemas de datos tabulares con relaciones no lineales.

A pesar de su éxito estadístico, múltiples teóricos coinciden en una barrera comercial: la opacidad algorítmica. Para mitigarlo, investigadores en \textit{eXplainable Artificial Intelligence} (XAI) propusieron implementaciones \textit{post-hoc} agnósticas al modelo, siendo LIME y SHAP (SHapley Additive exPlanations) los de mayor adopción empírica por su robustez apoyada en la teoría de juegos cooperativos \cite{lundberg2017unified}. 

Simultáneamente, la perspectiva académica referente al Costo de Clasificación Errante destaca que en el entorno empresarial no todos los errores tienen el mismo peso. Penalizar falsos positivos (campañas de retención gastadas en vano) y falsos negativos (el ingreso de vida del cliente o \textit{LTV} perdido) demanda enfoques \textit{Profit-driven} \cite{verbeke2012new}. A pesar de estos avances aislados, existe una notoria brecha en la literatura contemporánea respecto a estudios aplicados que operen las predicciones y sus desgloses causales XAI dentro de sistemas interactivos de Inteligencia de Negocios (BI) para optimización financiera en tiempo real.


\section{Metodología Estratégica Computacional}
En esta sección delineamos el flujo de arquitectura computacional híbrida (Machine Learning y Business Intelligence) diseñada para identificar, medir y accionar el desgaste de clientes en el entorno de telecomunicaciones.

\subsection{Procesamiento de Datos y Componente Predictivo}
El punto de partida del diseño consiste en el modelo de regresión o ensembilado. Asumimos un vector de covariables $x_i$ que engloba tanto métricas demográficas como el uso mensual y ciclo de servicio para un cliente determinado $i$. La variable dependiente del interés es la fuga binaria del cliente observada de forma histórica: 

\begin{equation}
y_i \in \{0, 1\}
\end{equation}

Empleamos el algoritmo Extreme Gradient Boosting (XGBoost) para optimizar un objetivo de pérdida logarítmica (log-loss), en vista de su capacidad excepcional para capturar no linealidades en datos tabulares y ser ávidamente decodificado a través del formalismo de valores de Shapley. 
La probabilidad inferida por nuestro modelo resultante se denotará de aquí en adelante como el puntaje de predicción probabilístico (Score Predictivo): $\hat{y}_i = P(y_i=1 | x_i)$.

\subsection{Inteligencia Artificial Explicable (XAI) mediante Valores SHAP}
Para transicionar de un ambiente de "caja negra" a un Sistema de Soporte de Decisiones (DSS), integramos SHAP (SHapley Additive exPlanations). Los valores SHAP transforman la salida predictiva $\hat{y}_i$ en combinaciones lineares aditivas del aporte causal de cada atributo individual:

\begin{equation}
\hat{y}_i = \phi_0 + \sum_{j=1}^{M} \phi_{ij}
\end{equation}

donde $\phi_0$ representa el valor base esperado del dataset y $\phi_{ij}$ es el valor cuasi-causal aportado por la variable de servicio $j$. Esta arquitectura es fundamental, pues en el sistema final (BI) los agentes comerciales pueden dirigir retenciones leyendo exclusivamente los vectores de peso $\phi$ locales del cliente sin reentrenar el motor principal.

\subsection{Arquitectura Interactiva Basada en Rentabilidad (Profit-driven BI)}
Tradicionalmente, en la literatura analítica, la precisión se evalúa con el AUC o la métrica F1. Sin embargo, para responder a nuestra hipótesis, proponemos una \textit{Matriz de Confusión Guiada por Rentabilidad} implementada virtualmente. 
Se define el retorno de acción (ROI) basado en costos asimétricos:
\begin{itemize}
    \item \textbf{Costo de Retención ($C_R$):} El incentivo económico enviado bajo el criterio de que el usuario fugara (usualmente falsos positivos $FP$).
    \item \textbf{Ingreso Salvado ($I_S$ o $LTV$):} Valor conservado cuando un cliente efectivamente en riesgo acepta la retención (verdaderos positivos $TP$).
\end{itemize}

El Dashboard de Streamlit (aplicación base BI) realiza el cruce instantáneo entre los scores del modelo algorítmico y los valores estáticos financieros introducidos interactivamente por los directivos. El gerente define un Umbral de Decisión Acotado ($\tau$) que transforma probabilísticamente al sujeto logrando estimar la rentabilidad macro por cohorte evaluada: 

\begin{equation}
  \text{Ganancia Neta} = \left( \sum_{i=1}^{N_{\text{TP}}} (I_S - C_R) \right) - \left( \sum_{i=1}^{N_{\text{FP}}} C_R \right)
\end{equation}

Esta orquestación garantiza que el modelo predictivo posea una injerencia mensurable en el ecosistema comercial de la empresa. En el siguiente escenario (sección empírica), expondremos la validación de estos pilares utilizando nuestra integración con un set empírico Telco.


\section{Estrategia Empírica y Modelo}
En este apartado detallamos las propiedades del entorno experimental, el conjunto de datos administrado y la estrategia algorítmica utilizada para evaluar el framework teórico planteado en la sección anterior.

\subsection{Conjunto de Datos Financieros}
Nuestra evaluación empírica descansa sobre el \textit{IBM Telco Customer Churn Dataset}, un estándar en la literatura analítica conformado por registros agregados referenciados en un eje de tiempo transversal de 7,043 clientes. 

Las dimensiones cardinales del conjunto de datos incluyen:
\begin{itemize}
    \item \textbf{Atributos demográficos:} Incluyen perfiles como titularidad, dependientes y estado civil.
    \item \textbf{Suscripción a Servicios:} Variables categóricas denotando seguridad en red, protección de dispositivo y volumen contratado.
    \item \textbf{Factor Financiero:} \textit{MonthlyCharges} (Cargos Mensuales) y \textit{TotalCharges} (Cargos Acumulados), los cuales servirán paramétricamente para extraer el proxy del LTV (Valor de Vida del Cliente).
\end{itemize}

Durante el pre-procesamiento, vectorizamos todas las características categóricas (One-Hot / Labeling). Los clientes con valores nulos inicializados en cargos acumulados se reclasificaron con valor 0 (clientes nuevos en su primer mes). La matriz purgada constó del 100\% de las trazas iniciales para el ensamble.

\subsection{Diseño Experimental y Modelamiento Predictivo}
Para la división transversal, aplicamos una segmentación estratificada reteniendo el 80\% ($N=5,634$) de los perfiles en el conjunto de entrenamiento (Train) y reservando el 20\% ciego ($N=1,409$) para inferencia causal y test empírico. 

El núcleo predictivo consta de un ensamble de \textit{Gradient Boosting Trees} (XGBoost). Este se configuró con profundidad restrictiva (\textit{max\_depth} = 5) y tasa de aprendizaje conservadora (\textit{learning\_rate} = 0.1) acoplado a un bloque de 100 estimadores. Las hiperparametrizaciones minimizan deliberadamente eventos de \textit{overfitting} sobre el espacio de características nominales de 20 variables.

\subsection{Proyección a Inteligencia de Negocios y Función de Costos}
Una vez obtenidas las proyecciones en el ambiente de inferencia (Test Set), transferimos los perfiles a nuestra plataforma virtual. En lugar de limitarnos a reportar la Curva Característica Operativa (ROC), evaluamos el alcance monetario sometiendo todo el lote de pruebas ($N_{test}$) a una política comercial hipotética de rescate.

Se configura a priori un costo operativo $C_R=50$ USD referenciando un descuento asimilable por el sistema e ingresado interactiva, con un LTV de retención proyectado $I_S=300$ USD. El dashboard (BI) reescala entonces el umbral crítico $\tau$ (Threshold) para dictaminar empíricamente qué porcentaje de alertas modeladas por XGBoost resultan factibles y eficientes sin decantar en un gasto infructuoso de adquisición. Simultáneamente, por cada perfil, se conjuga la métrica de SHAP de manera tabular para explicarle a los visores de Inteligencia de Negocios los impulsores demográficos detrás de su propensión al desgaste.


\section{Resultados y Simulaciones en Tableros BI}
La implementación computacional validó con éxito el modelo XGBoost en la sección transversal de prueba de IBM Telco ($N=1,409$). Desde el punto de vista algorítmico estándar, el clasificador capturó correctamente las señales multivariadas de desgaste reportando un rendimiento elevado, lo que fundamenta sólidas capacidades deductivas como motor estadístico principal.

Sin embargo, el eje central de nuestros resultados emerge al acoplar estas predicciones estáticas en el motor dinámico de Inteligencia de Negocios (BI). Durante la simulación de intervención interactiva ($C_R = 50$ USD como costo de retención e $I_S = 300$ USD de rescate), identificamos que optimizar la curva de decisión empujando un umbral ($\tau$) conservador sobre los perfiles propensos previene el desgaste comercial de campañas masivas ciegas. El cruce entre el número de Falsos Positivos no intervenidos (eludiendo la purga por $C_R$) y de Verdaderos Positivos reconvertidos, dictaminó una rentabilidad neta teórica escalable bajo la matriz de confusión interactiva.

Adicionalmente, el volcado gráfico de los valores de SHAP locales proveyó a la interfaz BI de una anatomía del riesgo unívoca. Esto facultó operativamente a los tomadores de decisión para diagnosticar a los clientes evaluados de forma unitaria, separando inmediatamente entre aquellos individuos empujados al abandono por factores financieros directos (\textit{MonthlyCharges}) en contraposición a aquellos sin lazos de lealtad estructurales técnicos (ausencia de \textit{TechSupport} o \textit{OnlineSecurity}).


\section{Conclusión}
Hemos demostrado a través del presente documento empírico que la maximización del valor de inteligencia predictiva no rige de forma exclusiva en la precisión estadística, sino fundamentalmente en su nivel de explicabilidad y acoplamiento directo con la estructura de costos empresariales reales. 

El desarrollo guiado longitudinalmente estableció una arquitectura fluida donde los datos transaccionan desde sus orígenes en bruto, se destilan computacionalmente mediante potentes ensambles de árboles con decodificación humana XAI (SHAP), y culminan de forma intuitiva como tableros interactivos para los niveles directivos. Este trabajo ha trasladado exitosamente los tradicionales prototipos puros orientados a la academia e integró las perspectivas de Inteligencia de Negocios para prescribir rentabilidad.

Estudios subsecuentes deberían apuntar al levantamiento de diseños causales experimentales (A/B testing iterativos controlados) que comprueben que las palancas expuestas como factores causadores de la probabilidad final por los algoritmos reaccionen, una vez operadas activamente por los operadores de servicio, verdaderamente mitigando la huida prospectiva.


\bibliographystyle{plainnat}
\bibliography{../Bibliography_base}

\end{document}
